% Reader-ready C05 evidence continuation for ED26-BM-RAT2021-15-1.
% Root-owned integration: these records extend the semantic/divergent-thinking
% section; core chapter and original intakes were intentionally not edited.
\subsection{The preschool experiment published in 2020: an author-reproduced record}
\label{sec:p19-r151-2020-author-record}

The 2020 article by de Chantal, Gagnon-St-Pierre and Markovits is a different
study from the two 2017 cohorts. Its complete author-reproduced main text is
Chapter~VIII, pp.~95--134, of de Chantal's August 2021 UQAM thesis.
The thesis's own bibliography identifies the published article as
\emph{Divergent Thinking Promotes Deductive Reasoning in Preschoolers},
\emph{Child Development} 91(4), 1081--1097, DOI 10.1111/cdev.13278.
The following methods and results are attributed to that exact thesis
presentation, not asserted byte-identical to the unavailable journal PDF.
The reproduced text refers to an appendix of full reasoning stimuli and
category coding examples that is not included with this study in the thesis;
those missing protocol details are not borrowed from its separate 2017 appendix
\cite[pp.~92--134,151]{DeChantal2021ThesisSelectedStudy}.

The experiment tested whether a preceding generation exercise helps children
distinguish a necessary inference from an invited but uncertain converse,
and whether explicit examples alter that benefit. Its comparison was between
generation without examples (DT-free), generation with examples (DT-examples),
and a preference-report control. This addresses the lack of a neutral
comparison in the 2017 generation-versus-card-sort study; it does not create
a longitudinal developmental test.

One hundred twenty children were recruited through eight day-care centers
in the Montreal region. The reported age range was 46--79 months, mean
60.0 months, SD~5.94; there were 65 boys and 55 girls. That range extends
outside the abstract's nominal four-to-five-year description. The sample
was predominantly Caucasian and drawn from middle-class neighborhoods.
Parental consent confirmed that French was mainly spoken at home; three
children also spoke another language and were reported to complete the
experiment without difficulty. This is a restricted local recruitment frame,
not a representative sample of all preschool children
\cite[pp.~107--112]{DeChantal2021ThesisSelectedStudy}.

Four experimenters received common training and were blind to the hypothesis.
Individual sessions used an animated, recorded presentation in a separate
day-care room. The experimenter sat beside the child, wrote responses,
handled questions and restarted the video after each required answer. A
white board limited surrounding distractions and possible environmental cues.
Children were randomly assigned one of the three preliminary activities,
then received the reasoning task and finally a knowledge/comprehension check.
Sessions lasted seven to twelve minutes and ended with a sticker. Hypothesis
blinding does not establish condition blinding: the example manipulation was
part of the presented activity.

Generation prompts asked for ways of making noise, giving gifts and getting
dirty, always in that order. DT-examples supplied two examples for each
prompt: a whistle/paper tearing, a soft toy/chocolate, and mud/ketchup.
DT-free used the same themes but requests for more possibilities without
those examples. The control instead asked about favorite colors, games and
foods, also in fixed order, to elicit roughly comparable verbal output.
Experimenters could not give feedback or encouragement but could prompt for
different ideas after a pause of approximately five seconds, and moved on
only when the child reported having no more ideas. An exercise without
examples was therefore not an absence of instruction or interaction
\cite[pp.~108--111]{DeChantal2021ThesisSelectedStudy}.

The reasoning task had four familiar category--property premises, formulated
as universal statements. Each received MP first, then AC, with standardized
requests for explanations. This reverses the AC-before-MP sequence in the
2017 record. The main text's examples and justifications concern dogs/legs,
trucks/wheels, flies/wings and oranges/edibility; they do not supply the
missing appendix's complete wording. The subsequent check explicitly asked
for an alternative to each premise's category, stopping after one answer.
Passing such a prompted retrieval check does not demonstrate spontaneous
retrieval during the earlier inference task.

Three children had incomplete sessions, two departed from the protocol,
five were excluded following educators' reports of verbal production or
comprehension difficulties, and seven failed the four check questions.
The latter exclusions were four from DT-examples, one from DT-free and
two from control. The analyzed cohort was therefore 103, not 120:
51 boys and 52 girls, age 46--79 months, mean~60.3, SD~5.69.
Final group sizes were 32, 32 and 39 for DT-examples, DT-free and control,
respectively. Their mean ages were 59.3, 60.2 and 61.1 months, with SDs
5.60, 5.96 and 5.55. The main text does not give each initially randomized
group's size or a full condition breakdown of the other ten exclusions;
equal initial allocation is not invented. The authors report a similar
pattern when the seven check failures were included, but do not supply
the corresponding numerical estimates here
\cite[p.~112]{DeChantal2021ThesisSelectedStudy}.

In this record, \(p\) denotes a statistical p-value, \(F\) an analysis-of-variance
statistic and \(\chi^2\) a reported chi-square statistic, not a logical
antecedent, connective or representation model. A correct MP answer expressed
certainty; a correct AC answer expressly denied certainty. ``I don't know''
was scored incorrect for both. Each separate inference score ranged from
zero to four. The paired deductive score awarded a point only when both
answers for the same premise were correct, also ranging from zero to four.
Overall MP and AC means were 3.87, SD~.50, and 1.25, SD~1.61;
paired-score mean was 1.18, SD~1.56, with mode zero. Unlike the 2008
categorical task below, a nonspecific uncertainty response was not grouped
with rejection. Thus changing score rules can change apparent competence.

An age-adjusted Poisson loglinear analysis reported a condition effect,
Wald \(\chi^2(2)=13.63,p<.01\); its age term was
\(\chi^2(1)=2.79,p=.10\). Paired-score means were 1.97, SD~1.75,
for DT-free, 1.09, SD~1.47, for DT-examples, and .62, SD~1.21,
for control. Bonferroni comparisons gave DT-free versus control
\(p<.001\), DT-free versus examples \(p<.05\), and examples versus
control \(p=.06\). These are mean counts out of four premise pairs,
not percentages of children with a domain-general competence. Figure~8.1's
error bars are standard errors, not confidence intervals for those contrasts
\cite[pp.~112--114]{DeChantal2021ThesisSelectedStudy}.

The source adds a log-score ANCOVA, reporting
\(F(2,99)=7.99,p<.01,\eta_p^2=.139\). It does not specify how zeros
were handled before a logarithm; no offset or exclusion rule is supplied here.
Its age-term line literally prints \texttt{p = 12}, which is not a valid
p-value and is not silently repaired. A nonsignificant goodness-of-fit
comparison does not establish that a bounded zero-to-four outcome has an
exact Poisson law. These reporting and model boundaries limit exact
reproduction without invalidating the distinction among the observed group
means; the Poisson and ANCOVA procedures are analyses of the same data,
not two independent replications.

The generation scores also need their operational definitions. Redundant
ideas and literal repetitions of supplied examples were removed. Labels
were standardized to avoid counting paraphrases as different ideas. The
two DT conditions were pooled to construct one frequency lexicon per theme,
yielding 93, 133 and 90 distinct ideas. Research assistants formed thematic
categories, and two coders assigned ideas to them, resolving disagreements
after Cohen kappa values .77, .75 and .64. Flexibility counted represented
categories; each uncategorized idea received its own flexibility point.
Originality counted ideas occurring once in that pooled sample divided by
the child's fluency, then averaged over the three exercises. This is
sample-relative statistical rarity, not an invariant property of an idea.
The main text does not state the division convention for an exercise with
zero ideas or reproduce the promised category appendix
\cite[pp.~116--119]{DeChantal2021ThesisSelectedStudy}.

Table~8.2 prints fluency, flexibility and originality means of
3.13, 1.98 and .19 in DT-examples, versus 3.45, 2.05 and .30 in
DT-free; corresponding SDs are 2.21, 1.43 and .21 versus 1.72,
.81 and .17. The earlier verbal-output paragraph instead prints 3.12
and 3.47 for the DT fluency means, and 3.23 for control. Those small
reporting differences are retained rather than declared identical.
The three-condition verbal-output analysis had \(F(2,99)=.23,p=.80\);
nonsignificance is not proof of equal verbal activation. In the DT-only
age-adjusted analyses, originality differed,
\(F(1,61)=5.49,p<.05,\eta_p^2=.083\), whereas fluency and flexibility
comparisons had \(p=.21\) and \(p=.38\). Reported Cronbach alpha across
three exercises was .74 for fluency, .71 for flexibility and only .40 for
the originality ratio; the undivided unique-idea count had alpha .67.
The lower ratio reliability matters to individual-difference claims
\cite[pp.~116--120, table~8.2]{DeChantal2021ThesisSelectedStudy}.

Age-controlled correlations of reasoning with fluency, flexibility and
originality were .45, .48 and .36 in DT-examples, versus .15, .10 and
.02 in DT-free. The reported between-condition correlation comparisons
did not pass .05: their p-values were .10, .054 and .08, respectively.
Significance in one group but not another does not itself establish different
population correlations. The proposed attenuation of individual differences
by DT-free is an interpretation, not demonstrated causal mediation by
originality. No interval estimates for these correlations or the principal
group contrasts are printed
\cite[pp.~120--126, table~8.3]{DeChantal2021ThesisSelectedStudy}.

The explanation data provide a further but dependent observation. Alternative
antecedents occurred in 55.0 percent of explanations following correct AC
answers versus 7.6 percent following incorrect ones. Disablers occurred in
83.3 percent of explanations following the relatively rare incorrect MP
answers. The reported justification-analysis denominators, 392 and 408,
count responses from the 103 children, not new participant samples.
Explanations elicited after judgments do not uniquely identify the process
that produced the judgment, and the study was not designed to identify
condition-specific qualitative changes in those explanations
\cite[pp.~114--116,125]{DeChantal2021ThesisSelectedStudy}.

Randomized preliminary activities support a task-specific short-term
comparison after the stated exclusions. The stronger no-example result
constrains a mere-more-answers explanation, but does not isolate originality
from every difference in semantic cueing, attention, motivation or task frame.
The authors explicitly leave duration, broader individual differences and
socioeconomic generalization open. Nothing here establishes a stable age
equivalence, transfer to abstract or causal premises, enduring improvement,
or a universally effective educational practice. The absent original appendix
and journal-version comparison remain material limits on exact protocol
qualification, not reasons to substitute the abstract for this evidence.

\subsection{The 2008 contrast: necessity, likelihood and confidence}
\label{sec:p19-r151-2008-record}

Markovits and Thompson's three original studies separate three questions:
whether a conclusion is necessary, how likely it is, and how confident the
reasoner is in the likelihood judgment. Their categorical answers and
five-point ratings are different measures. The paper's proposed two-strategy
interpretation is a psychological hypothesis, not a consequence of having
two response formats
\cite[pp.~1066--1069]{MarkovitsThompson2008Development}.

Study~1 involved 56 children from rural state schools in southwest England:
17 in Level~1, mean age 6 years 5 months, 2 girls/15 boys;
13 in Level~2, mean 7 years 1 month, 5 girls/8 boys; and 26 in
Level~4, mean 8 years 7 months, 13 girls/13 boys. The paper calls these
nominal six-, seven- and nine-year cohorts; it does not report individual
age ranges, a probability recruitment frame or response rates. Four
videotaped sequences varied premise order and whether categorical or
probabilistic questions came first, systematically from child to child,
not by a stated random allocation. Individual sessions lasted about ten
minutes. MP always preceded AC for a given premise.

Premises concerned cars/motors and hedgehogs/four legs, chosen to provide
accessible alternatives and few disablers. Categorical instructions asked
for certainty assuming the premise true, followed by a justification.
``Yes'' counted as acceptance; both ``no'' and ``not sure'' counted as
nonacceptance. The likelihood task used five increasing red bars, labeled
from not sure to very, very sure. Practice displays contained one, five
or nine red trees out of ten, with repeat instruction when needed.
The article supplies English formulations; it does not separately state
native-language eligibility for these English-school cohorts
\cite[pp.~1069--1071]{MarkovitsThompson2008Development}.

\begin{table}[htbp]
\centering\small
\begin{tabular}{lrrrr}
\hline
Level & Pupils & MP accept (\%) & AC accept (\%) & AC rating\\
\hline
1 & 17 & 88.2 & 61.8 & 3.57\\
2 & 13 & 88.5 & 11.5 & 3.38\\
4 & 26 & 84.6 & 25.0 & 2.72\\
\hline
\end{tabular}
\caption{Study~1's categorical acceptance percentages and separate AC
likelihood ratings on a one-to-five scale, combined over two premises.
The likelihood numbers are not numerical probabilities or correctness scores.}
\label{tab:p19-r151-2008-study1}
\end{table}

Table~\ref{tab:p19-r151-2008-study1} shows that categorical AC acceptance
does not form a strictly monotone age ladder in these small cohorts.
Mean MP likelihood ratings were 4.56, 4.50 and 4.56, respectively.
An instruction-task comparison therefore need not imply that all of one
cohort's inferences are superior. Counterexample explanations were coded
by two experimenters with fewer than five percent initial disagreements,
resolved by discussion. Their correlations with categorical AC acceptance
were -.69, -.61 and -.75 across the three levels; with AC likelihood
ratings they were -.56, -.59 and -.01. The paper reports a within-oldest-
cohort difference between those correlations. These relationships involve
elicited explanations, not randomized counterexample use, and the paper
prints no confidence intervals for them
\cite[pp.~1071--1073, tables~1--3]{MarkovitsThompson2008Development}.

Study~2 used another 53 pupils from the same kind of southwest-English
schools: 21 in Level~1, mean 6 years 5 months, 11 girls/10 boys;
18 in Level~2, mean 7 years 7 months, 8 girls/10 boys; and 14 in
Level~4, mean 9 years 8 months, 8 girls/6 boys. The visible high-likelihood
array contained four red squares, one red circle and four blue circles;
the low-likelihood array contained one red triangle, four red circles and
four blue circles. The rules linked squares or triangles to red.
Conditional on a red shape and equal selection among those red shapes,
the target proportions are four-fifths and one-fifth. Both arrays still
contain a red counterexample to the converse, so neither AC conclusion
is deductively necessary. Actual selection probabilities require that
equal-selection assumption, not just counting the picture.

The procedure repeated the four systematically varied sequence formats
but omitted justification requests. Average AC likelihood ratings were
3.87 for the high-likelihood array versus 2.56 for the low-likelihood
array, \(F(1,47)=22.64,p<.001\). This supports discrimination, not
numerical calibration of a probability distribution. Categorical AC
acceptance averaged 49.1 versus 28.3 percent by array. It also differed
with task order: 25.9 percent when categorical questions came first,
versus 51.9 percent when likelihood questions came first. The reported
order comparison was \(\chi^2(1)=6.34,p<.01\). Likelihood ratings
showed no reported order effect; that nonsignificance does not prove exact
order invariance. Grade-specific order cells contained only seven to
eleven children, and the stimuli/instructions differ from the preschool
generation experiment
\cite[pp.~1073--1075, table~4]{MarkovitsThompson2008Development}.

The printed probabilistic example on p.~1074 repeats the triangle
antecedent in its question instead of asking about the consequent, unlike
the surrounding inference description. Source-image inspection confirms
this wording discrepancy; it is not repaired by inventing a verified
administration script. The record preserves the stated paradigm and
published results while leaving that exact example's protocol standing
qualified.

Study~3 tested a separate 26 native French-speaking Grade~2 children in
a Brussels public elementary school, mean age 7 years 3 months,
13 girls and 13 boys. Two drawings showed eight red squares and one
blue square, or eight circles and one triangle. Each child first rated
the likelihood of choosing the frequent kind with eyes closed, then
confidence in that rating; the rare-kind questions followed. The same
five-point scale was used for both judgments. Drawing order varied
between two halves of the cohort; frequency-question order did not.
Sessions lasted approximately five minutes.

For the square display, likelihood means for frequent/rare outcomes were
3.61/2.62, whereas corresponding confidence means were 3.81/3.92.
For circles/triangle, the values were 3.54/2.88 and 3.69/3.96.
A low-likelihood event can therefore receive a relatively confident
rating without being rated likely. The rare-event likelihood/confidence
comparisons were \(t(25)=3.45,p<.001\) and
\(t(25)=3.11,p<.01\). The main likelihood ANOVA prints
\(F(1,23)=7.11,p<.01\); its denominator is retained rather than
used to invent an unreported exclusion record. The article gives no
intervals for the displayed means. This is evidence for a distinction
within this instructed rating task, not a demonstration that every young
child distinguishes confidence from probability in every context
\cite[pp.~1075--1076, table~5]{MarkovitsThompson2008Development}.

The article itself qualifies a simple analytic/heuristic dichotomy:
retrieving an alternative can be automatic when strongly associated,
but effortful when weakly associated; explicit representation is another
process. Its cross-sectional task differences do not measure maturation
of one pure system. Comparing these records with preschool-generation
scores also requires different recruitment, premise content, uncertainty
coding and score denominators to remain visible. Neither the 2008
instruction contrast nor the later generation prime establishes an
age-stage equivalence or a universally optimal reasoning strategy
\cite[pp.~1076--1078]{MarkovitsThompson2008Development}.

